% =============================================================================
%  Clinzo --- Walkthrough Video: full verbatim narration script
%
%  Overleaf-ready. Standard TeX Live packages only, compiles with pdfLaTeX,
%  no shell-escape required. Upload this single .tex file and Recompile.
% =============================================================================

\documentclass[11pt,a4paper]{article}

\usepackage[T1]{fontenc}
\usepackage[utf8]{inputenc}
\usepackage[margin=2.1cm]{geometry}
\usepackage{xcolor}
\usepackage{booktabs}
\usepackage{longtable}
\usepackage{array}
\usepackage{enumitem}
\usepackage{listings}
\usepackage{tcolorbox}
\usepackage{fancyhdr}
\usepackage{titlesec}
\usepackage[hidelinks]{hyperref}

% --- palette ---
\definecolor{ink}{HTML}{1B2733}
\definecolor{accent}{HTML}{1D6FB8}
\definecolor{soft}{HTML}{5B7085}
\definecolor{shell}{HTML}{F4F6F8}
\definecolor{rule}{HTML}{D8E0E8}
\definecolor{stagec}{HTML}{8A4B08}

\color{ink}

\titleformat{\section}{\Large\bfseries\color{accent}}{\thesection}{0.6em}{}
\titleformat{\subsection}{\large\bfseries}{\thesubsection}{0.6em}{}

\pagestyle{fancy}
\fancyhf{}
\lhead{\small\color{soft}Clinzo --- Walkthrough Narration Script}
\rhead{\small\color{soft}\thepage}
\renewcommand{\headrulewidth}{0.4pt}

\lstdefinestyle{sh}{
  basicstyle=\ttfamily\small,
  backgroundcolor=\color{shell},
  frame=single,
  rulecolor=\color{rule},
  breaklines=true,
  columns=fullflexible,
  keepspaces=true,
  showstringspaces=false,
  framexleftmargin=6pt,
  framexrightmargin=6pt,
  aboveskip=8pt,
  belowskip=8pt,
}
\lstset{style=sh}

% --- helpers ---

% Verbatim narration --- everything inside is read aloud, word for word.
% Paragraphs are spaced rather than indented: this is read aloud off a screen,
% so the visual gap between beats is what stops you losing your place.
\newtcolorbox{script}[1][]{
  colback=white, colframe=accent, boxrule=1.1pt,
  left=9pt, right=9pt, top=7pt, bottom=7pt, arc=2pt,
  fonttitle=\bfseries\small, title={SAY THIS},
  before upper={\setlength{\parskip}{7pt}\setlength{\parindent}{0pt}},
  #1}

% Stage direction --- do this, do not read it out.
\newcommand{\stage}[1]{%
  \par\vspace{4pt}\noindent
  \textcolor{stagec}{\textbf{[\,DO:}\ \textit{#1}\textbf{\,]}}\par\vspace{4pt}}

\newtcolorbox{note}[1][]{colback=shell,colframe=rule,boxrule=0.7pt,
  left=7pt,right=7pt,top=5pt,bottom=5pt,arc=2pt,#1}

\newcommand{\seg}[4]{%
  \vspace{10pt}
  \subsection*{#1}
  \noindent{\color{soft}\small Runs #2\quad$\cdot$\quad#3 words\quad$\cdot$\quad#4}
  \par\vspace{4pt}\hrule\vspace{6pt}}

% =============================================================================
\begin{document}

\begin{center}
  {\LARGE\bfseries Clinzo --- Doctor Slot Scheduling}\\[5pt]
  {\Large\color{soft} Walkthrough video: full narration script}\\[10pt]
  {\small Total spoken length \textbf{$\approx$12 min 45 s} at a normal pace
   (hard limit 15).\\
   Repository: \texttt{github.com/Subharjun/clinzo}}
\end{center}

\vspace{8pt}
\hrule
\vspace{12pt}

\begin{note}
\textbf{How to use this document.} Everything inside a blue \textsc{say this}
box is written to be read aloud exactly as printed. It is already paced ---
roughly 145 words per minute, which is a relaxed speaking speed. The orange
\textbf{[\,DO:\,]} lines are stage directions: perform them, do not read them.

Read it through twice before recording. You do not need to memorise it --- put
it on a second monitor or print it. Small deviations are fine; just do not
speed up to fit, because rushing the concurrency section is the one thing that
would actually damage the submission.
\end{note}

% =============================================================================
\section{Before you hit record}

\subsection*{1. Start clean}

\begin{lstlisting}
cd ~/Desktop/Clinzo
docker compose up -d postgres redis   # wait until both report healthy
npm run seed                          # accounts, availability, slots (idempotent)
npm run demo:reset                    # clear bookings/holds, all slots AVAILABLE
npm run dev                           # API on :3000
\end{lstlisting}

\begin{note}
\textbf{Run \texttt{demo:reset} immediately before recording.} The opening shot
depends on Dr Mehta's first three slots being free, because they are the worked
example from the brief: \textbf{10:00, 10:20, 10:40}. If you have been clicking
around, those will be booked and the strongest visual in the video is gone.

\textbf{\texttt{npm run seed} alone will not fix that.} The seed is idempotent
by upsert --- deliberately, so it is safe to re-run against a database you are
working in --- which means it never deletes bookings or holds.
\texttt{npm run demo:reset} is the one that clears the transactional state and
returns every slot to \texttt{AVAILABLE}.
\end{note}

\subsection*{2. Windows you will need}

\begin{enumerate}[leftmargin=*,itemsep=2pt]
  \item \textbf{Browser tab 1} --- demo console, \texttt{http://localhost:3000/app}
  \item \textbf{Browser tab 2} --- API docs, \texttt{http://localhost:3000/docs}
  \item \textbf{Terminal} --- dark theme, font at 16--18pt
  \item \textbf{Editor} --- with these four files already open as tabs:
    \begin{itemize}[leftmargin=1.2em,topsep=2pt,itemsep=1pt]
      \item \texttt{src/services/slot-generator.ts}
      \item \texttt{prisma/migrations/20260726024738\_init/migration.sql}
      \item \texttt{src/services/availability.service.ts}
      \item \texttt{src/repositories/outbox.repository.ts}
    \end{itemize}
\end{enumerate}

\subsection*{3. Recording checklist}

\begin{itemize}[leftmargin=*,itemsep=2pt]
  \item 1920$\times$1080. Hide the bookmarks bar. Close unrelated tabs.
  \item Silence notifications --- a Slack popup mid-take costs you a re-record.
  \item Test the microphone first. Bad audio sinks a good demo.
  \item Record \textbf{one take per segment}. Re-recording ninety seconds is
        painless; re-recording thirteen minutes is not.
  \item Reset between takes with \texttt{npm run demo:reset} (about a second).
        Then hard-reload the console --- the slot listing is cached for 15
        seconds, so a soft refresh can still show the old state.
\end{itemize}

% =============================================================================
\section{Shot list}

\begin{longtable}{@{}p{0.7cm} p{6.1cm} p{1.5cm} p{1.5cm} p{4.4cm}@{}}
\toprule
\textbf{\#} & \textbf{Segment} & \textbf{Length} & \textbf{Ends} & \textbf{Window} \\
\midrule
\endhead
1 & The problem and the guarantee   & 0:52 & 0:52  & Console \\
2 & Stack and architecture          & 1:19 & 2:11  & Editor \\
3 & Slot generation and timezones   & 2:16 & 4:27  & Console \\
4 & Booking lifecycle               & 1:52 & 6:19  & Console \\
5 & \textbf{The concurrency proof}  & 2:46 & 9:05  & Console, terminal, SQL \\
6 & Retroactive availability change & 1:12 & 10:17 & Editor \\
7 & Events, outbox and workers      & 1:08 & 11:25 & Editor \\
8 & Ops surface and CI              & 1:02 & 12:27 & Browser, terminal \\
9 & Trade-offs and close            & 0:45 & 13:12 & Console or camera \\
\bottomrule
\end{longtable}

% =============================================================================
\section{The narration}

% ---
\seg{Segment 1 --- The problem and the guarantee}{0:52}{125}{Demo console, Dr Mehta's slots on screen}

\stage{Start on the demo console with slots visible. Do not click anything yet.}

\begin{script}
Hi, I'm Subharjun. This is Clinzo --- a doctor slot scheduling backend I built
for the backend engineering assessment.

The problem looks simple at first. A doctor says: I work Mondays, ten to six. A
patient needs to book one specific fifteen-minute appointment inside that
window. Turning a broad availability window into discrete bookable slots is the
easy half.

The hard half is this. Two patients want the same ten o'clock slot, and they
click at exactly the same moment. What happens?

So here is the one claim I want to make, and I'll spend the next twelve minutes
backing it up. This system will never double-book a slot. Not under load, not
across multiple servers, not ever. Let me show you why.
\end{script}

% ---
\seg{Segment 2 --- Stack and architecture}{1:19}{190}{Editor, repository tree}

\stage{Switch to the editor. Show the repository tree, then expand \texttt{src/}.}

\begin{script}
A quick tour of the stack before we get to the interesting part.

Node 22, TypeScript in strict mode, Express for HTTP, PostgreSQL through Prisma,
Redis for distributed locking, caching and rate limiting, and BullMQ for
background jobs. All of it containerised, with GitHub Actions running the full
suite on every push.

That's about ten thousand eight hundred lines of TypeScript, and roughly three
thousand more of tests.
\end{script}

\stage{Point at each folder as you name it: \texttt{routes/}, \texttt{controllers/}, \texttt{services/}, \texttt{repositories/}.}

\begin{script}
The architecture is conventionally layered, and I was strict about it.
Controllers do exactly three things: translate HTTP into service arguments,
call one service method, and translate the result back. No branching on domain
state, no database access. Every real decision lives in the service layer, and
everything that touches Postgres goes through a repository.
\end{script}

\stage{Open \texttt{src/services/slot-generator.ts}. Scroll slowly through the top of the file.}

\begin{script}
One module deserves a mention now, because it comes back later. The slot
generator is completely pure. No database, no clock, no input or output of any
kind --- availability goes in, slots come out. That's what makes the hardest
arithmetic in the system directly testable.
\end{script}

% ---
\seg{Segment 3 --- Slot generation and timezones}{2:16}{330}{Demo console, ``Find a slot''}

\stage{Switch to browser tab 1. Dr Anaya Mehta should already be selected.}

\begin{script}
Let's look at the running system. This is a small demo console I built on top
of the API --- plain HTML and JavaScript, served by the same Express app, no
framework and no build step. It exists purely to make the backend's behaviour
visible.

We're looking at Doctor Anaya Mehta, who works in Asia slash Kolkata. Fifteen
minute consultations, with a five minute buffer between them.

Now look at the first three slots. Ten o'clock. Ten twenty. Ten forty. Fifteen
minutes of appointment plus five minutes of buffer, so they start twenty
minutes apart. That is the worked example from the assessment brief, and it
falls straight out of that pure generator I just showed you.
\end{script}

\stage{Point at the three time rows inside a single slot card.}

\begin{script}
Here's the part I find more interesting. Every slot is stored in the database
as one UTC instant --- that's this row here. But it's rendered three ways: the
UTC instant, the viewer's local time, and the doctor's local time. And
critically, all three of those are computed by the server, not by the browser.
\end{script}

\stage{Change ``View times in'' to \textbf{America/New\_York}. Click \textbf{Search}.}

\begin{script}
Same slots. Same underlying instants. Now displayed in New York time, with the
doctor's Kolkata time alongside. Because when a patient in New York and a
doctor in India both talk about the ten forty appointment, they have to be
talking about the same moment.

Two bugs I hit here are worth mentioning, because they're the kind of thing
this domain is full of.

The first one: Luxon's plus-minutes function measures elapsed time. So on a
spring-forward day, midnight plus a hundred and eighty minutes is four in the
morning, not three. Every slot after the transition was silently shifting by an
hour. A test caught that one.

The second: I had a validator that happily accepted the string I-S-T. The
timezone database says that means Israel Standard Time. Almost everyone typing
it means India. That's a three and a half hour error on every single
appointment. So now I require the full Area slash Location form, and U-T-C is
the only bare name I allow.
\end{script}

% ---
\seg{Segment 4 --- Booking lifecycle}{1:52}{270}{Demo console, signed in as a patient}

\stage{Click \textbf{Sign in}, click the \textbf{Patient} chip, click \textbf{Sign in}.}

\begin{script}
Let me sign in as a patient and walk through the booking lifecycle.

Authentication is JWT, with short-lived access tokens and refresh token
rotation. Every refresh issues a new token and invalidates the old one. And if
a token is ever replayed --- which means it was stolen --- the entire session
family is revoked immediately.
\end{script}

\stage{Book the 10:00 slot. Let the network panel on the right update.}

\begin{script}
Booked, with a confirmation code. And notice this panel on the right. That's
every API call the page makes, with its status code, its timing, and its
request id. Every response carries that request id, so when a user reports a
problem, that single value takes you to the exact log line.
\end{script}

\stage{Click \textbf{Hold} on a different slot, then open ``Holds \& waitlist''.}

\begin{script}
Now let me hold a slot instead of booking it. A hold is the checkout window ---
a short reservation while a patient finishes paying. That countdown is real and
it's enforced by the server. When it expires, the slot returns to the pool
automatically. Nothing has to succeed for that to happen, and that matters,
because the failure path is the one that has to be reliable.
\end{script}

\stage{Open ``My bookings''. Click \textbf{Reschedule}. Wait for the toast, then click \textbf{Cancel}.}

\begin{script}
Over in my bookings, let me reschedule this appointment.

It moved to a new slot, and it kept the same booking id. Reschedule is one
atomic operation --- deliberately not a cancel followed by a rebook, because
that leaves a window where the patient has given up their appointment and
hasn't got the new one yet. Here, if the target slot is taken, the original
booking is still theirs.

And cancelling puts the slot straight back into public availability.
\end{script}

% ---
\seg{Segment 5 --- The concurrency proof}{2:46}{400}{Console, then terminal, then the migration SQL}

\begin{note}
\textbf{This is the segment that matters.} Three beats: show the race, prove
the lock is not doing the work, then show the constraint that is. If you are
short on time elsewhere, cut segments 6 through 8 --- never this one.
\end{note}

\subsubsection*{Beat 1 --- show the race}

\stage{Open the ``Concurrency lab'' view. Pause a moment on the three layer cards.}

\begin{script}
This is the core of the assessment, so I'm going to spend real time here.

Three layers guard a booking. A Redis distributed lock. A SELECT FOR UPDATE row
lock inside the transaction. And a partial unique index in Postgres. Only the
third one is the actual guarantee --- but let me show you the race first.
\end{script}

\stage{Leave contenders at 12. Click \textbf{Fire simultaneously}.}

\begin{script}
Twelve different patient accounts, all signed in, twelve booking requests for
one single slot, released simultaneously.

One two-oh-one Created. Eleven four-oh-nine Conflicts. About a hundred and
thirty milliseconds, end to end.

They're twelve separate accounts on purpose. The booking rate limiter allows
thirty requests per minute per user, so if I fired all twelve from one account,
I'd be showing you the rate limiter rather than the concurrency control.
\end{script}

\subsubsection*{Beat 2 --- prove the lock is not the guarantee}

\begin{script}
Now here's the claim I actually want to defend. That Redis lock is not what
makes this correct.

A lease-based distributed lock cannot be made safe. If a process pauses --- a
long garbage collection, say --- past the lock's expiry, it wakes up believing
it still holds a lock that Redis has already handed to somebody else. And then
it writes. Treating Redlock as a correctness mechanism is precisely how systems
double-book.

So in this system, the lock is a throughput optimisation. It stops a hundred
contenders from each holding a database connection while they queue on the same
row. It is not the guarantee.

And I can prove that, by removing it.
\end{script}

\stage{Switch to the terminal and run:}
\begin{lstlisting}
npx jest tests/integration/booking-concurrency.test.ts \
  -t "survives with the Redis lock removed"
\end{lstlisting}

\begin{script}
This test stubs the lock out completely --- no Redis coordination at all ---
fires the contenders, and still asserts exactly one winner.

There it is, passing. That test is the actual proof. Everything else is
performance.
\end{script}

\subsubsection*{Beat 3 --- show the constraint}

\stage{Open \texttt{prisma/migrations/20260726024738\_init/migration.sql} and scroll to line 437.}

\begin{lstlisting}
CREATE UNIQUE INDEX "bookings_one_confirmed_per_slot"
  ON bookings ("slotId") WHERE status = 'CONFIRMED';
\end{lstlisting}

\begin{script}
So what is holding the line? This, in the migration.

A partial unique index. Unique on slot id, where the status is confirmed. One
confirmed booking per slot.

Postgres will refuse the second write regardless of what the application code
believes, regardless of how many API replicas are running, and regardless of
whether Redis is even up.

And that WHERE clause is doing real work. It's scoped to confirmed bookings
only --- which is exactly why the cancellation you saw a minute ago frees the
slot up again. The cancelled row stops counting, so the slot can be re-booked,
and still only once.
\end{script}

\begin{note}
\textbf{Likely follow-up question, worth having ready:} why pessimistic locking
for bookings but optimistic locking for availability? Because booking
contention is concentrated by nature --- everyone wants the ten o'clock --- so
optimistic locking makes every loser burn a full transaction to discover
something a retry can never fix. Availability edits are the opposite: rare, and
a stale write there genuinely should be rejected. So those carry version
columns.
\end{note}

% ---
\seg{Segment 6 --- Retroactive availability changes}{1:12}{175}{Editor, availability.service.ts}

\stage{Open \texttt{src/services/availability.service.ts} and scroll to the \texttt{update()} header comment.}

\begin{script}
Here's the design question I found hardest, and I don't think there's a clean
answer to it.

A doctor publishes Monday, ten to six. A patient books three o'clock. Then the
doctor shrinks Monday to ten until two. There is now a confirmed appointment
sitting outside the doctor's stated availability.

Three options. Auto-cancel the booking --- that silently breaks a commitment
between two people, and the patient finds out later. Reject the availability
change outright --- now one appointment freezes the doctor's entire schedule.
Or keep the booking and report the conflict.

I chose the third. Future unbooked slots outside the new window get blocked.
Future booked slots are left completely untouched, and come back in the API
response as orphaned bookings. Past slots are never modified at all.

That does leave the system in an inconsistent state, and I want to be honest
about that. But it's a transient inconsistency, it's visible, and it's
actionable --- and no commitment between two humans got cancelled by a
background job.
\end{script}

% ---
\seg{Segment 7 --- Events, outbox and workers}{1:08}{165}{Editor, outbox.repository.ts}

\stage{Open \texttt{src/repositories/outbox.repository.ts}.}

\begin{script}
Briefly on the event pipeline, because there's one decision worth explaining.

When a booking is created, a notification has to go out. The naive
implementation writes the booking, commits, and then publishes the event. If
the process dies between those two steps, the appointment exists and the
patient is never told.

So instead, the event is written to an outbox table inside the same transaction
as the booking. Either both land or neither does. A separate relay then polls
that table using SKIP LOCKED, and publishes from there. That gives me
at-least-once delivery with idempotent consumers --- which is a far easier
problem to solve correctly than distributed atomicity.

One thing I want to be upfront about. The final delivery hop to an email
provider is stubbed --- it logs instead of sending. Everything around it is
real: recipient resolution, per-recipient timezone rendering, retries with
backoff, deduplication by event id. Swapping in a real provider is one function
body.
\end{script}

% ---
\seg{Segment 8 --- Ops surface and CI}{1:02}{150}{Browser tab 2, then terminal}

\stage{Switch to browser tab 2 (\texttt{/docs}). Scroll the endpoint list.}

\begin{script}
The operational surface, quickly.

Hand-written OpenAPI documentation covering the whole API. I wrote it by hand
rather than generating it, because generators describe types --- and what
actually matters to somebody integrating is behaviour. What a four-oh-nine
means on this endpoint is the useful part.
\end{script}

\stage{Visit \texttt{/health}, then \texttt{/metrics}. A couple of seconds each.}

\begin{script}
Health endpoints with real dependency checks, split into liveness and
readiness. Prometheus metrics. And the whole stack --- Postgres, Redis,
migrations, the API and the worker --- comes up with a single docker compose
command.

CI runs lint, formatting, typecheck, the full test suite against real Postgres
and real Redis, and a Docker build, on every push. There's also an explicit
check that the partial unique index still exists --- because if somebody
regenerated that migration and dropped it, every test would still pass and the
guarantee would be silently gone.
\end{script}

% ---
\seg{Segment 9 --- Trade-offs and close}{0:45}{110}{Console or camera}

\begin{script}
Two trade-offs to close on.

Slots are materialised as real rows, rather than computed on read. That costs
storage and it costs a generation job. But a computed slot has no identity ---
there's nothing for FOR UPDATE to lock, and nothing for a unique index to
protect. The guarantee would have had to live in application code, which is
exactly what fails under concurrency.

And on the daylight saving fall-back, a doctor loses about an hour of inventory
once a year. I took that over ambiguous local times.

Next steps would be partitioning these tables before they get large, calendar
sync, and a real notification provider.

Thanks for watching.
\end{script}

% =============================================================================
\section{If something goes wrong mid-recording}

\begin{longtable}{@{}p{5.2cm} p{9.8cm}@{}}
\toprule
\textbf{Symptom} & \textbf{Fix} \\
\midrule
\endhead
Slots don't start at 10:00 & \texttt{npm run demo:reset}, then hard-reload. \\
Health dot is red & \texttt{docker compose up -d postgres redis}, wait for healthy. \\
Race returns 429s & You re-ran it many times inside one minute. Wait 60 seconds. \\
Race returns more than one 201 & Stop recording. That would be a real defect --- do not narrate around it. \\
``This slot has already been booked'' & Left over from an earlier take. Run \texttt{npm run demo:reset}. \\
Reschedule fails & Widen the date range so a free slot exists to move into. \\
Console won't load & Check \texttt{npm run dev} is running, and that the URL ends in \texttt{/app}. \\
\bottomrule
\end{longtable}

% =============================================================================
\section{Quick reference card}

\subsection*{Accounts --- password \texttt{ClinzoDemo2026!}}

\begin{longtable}{@{}p{5.8cm} p{2.2cm} p{6.4cm}@{}}
\toprule
\textbf{Email} & \textbf{Role} & \textbf{Notes} \\
\midrule
\endhead
\texttt{dr.mehta@clinzo.test}     & Doctor  & Asia/Kolkata, 15 min + 5 min buffer \\
\texttt{dr.okafor@clinzo.test}    & Doctor  & America/New\_York, 30 min, no buffer \\
\texttt{patient1..12@clinzo.test} & Patient & Mixed timezones \\
\texttt{admin@clinzo.test}        & Admin   & UTC \\
\bottomrule
\end{longtable}

\subsection*{URLs}

\begin{longtable}{@{}p{5.8cm} p{8.6cm}@{}}
\toprule
\textbf{What} & \textbf{Where} \\
\midrule
\endhead
Demo console       & \texttt{http://localhost:3000/app} \\
API docs (Swagger) & \texttt{http://localhost:3000/docs} \\
Health             & \texttt{http://localhost:3000/health} \\
Metrics            & \texttt{http://localhost:3000/metrics} \\
\bottomrule
\end{longtable}

\subsection*{Numbers you may want to quote}

\begin{longtable}{@{}p{4.2cm} p{10.2cm}@{}}
\toprule
185                & tests passing (117 unit, 68 integration) \\
$\approx$10{,}800  & lines of TypeScript source \\
12                 & Prisma models \\
31                 & route registrations \\
1 and 11           & 201s and 409s in a 12-way race, in $\approx$130\,ms \\
\bottomrule
\end{longtable}

\vspace{10pt}
\begin{note}
\textbf{Before you upload.} Watch it back once, at normal speed, with audio on.
Check three things: the code on screen is legible, the race result is readable,
and the total is under fifteen minutes.
\end{note}

\end{document}
